% ============================================================================
% Chapter 5: AutoEvolve (Headers Only)
% ============================================================================

\newpage
\section{AutoEvolve}\label{sec:autoevolve}

\subsection{Beyond Pok\'{e}mon}\label{sec:beyond_pokemon}

\textit{[To be written. Position AutoEvolve as a domain-agnostic evolutionary loop requiring modular agent framework, sandboxing, and a global session log; leverages LLM-as-a-judge in a reset-free manner. AutoEvolve is inherently domain agnostic.]}

\subsection{Methodology}\label{sec:methodology}

\textit{[NOTE from Ch.~2 refactor: The following details were moved here from the Formalization and Related Work sections of Chapter~2, to be integrated into the methodology.]}

\textit{[RLM connection: The trajectory window in PokeAgent functions as an RLM-style workspace---the full history of the agent's execution (state observations, prompts, reasoning traces, actions) is maintained as a structured, programmatically accessible log that subagents process and summarize for the orchestrator. This is a departure from standard agent architectures that treat history as a flat, append-only sequence. AutoEvolve's mid-episode evolution can be understood as a form of recursive self-modification: the agent analyzes its own trajectory segments and uses this analysis to rewrite its harness components (system prompt, subagents, skills, memory), which in turn shape subsequent trajectory segments. Unlike the depth-$k$ recursion that RLMs formalize, AutoEvolve operates at depth-1.]}

\textit{[Harness spectrum: Define $\mathcal{H}_{\min}$ (minimalist), $\mathcal{H}_{\text{expert}}$ (human-designed), $\mathcal{H}_{\text{meta}}$ (meta-tool), $\mathcal{H}_{\text{auto}}$ (auto-evolved). Central question: whether $\mathcal{H}_{\text{auto}}$, evolved from $\mathcal{H}_{\min}$ through AutoEvolve's reset-free optimization, can approach the performance of $\mathcal{H}_{\text{expert}}$.]}

\textit{[In-context prompt evolution formalism: episodic setting (Eq.~for GEPA) vs. reset-free setting (Eq.~for AutoEvolve). Four-pass structure over prompts, subagents, skills, and memory.]}

\subsubsection{Algorithm}\label{sec:algorithm}

\textit{[To be written. Four-pass evolutionary optimization: system prompt, subagent registry, skill library, memory store. Scheduling (warmup steps, evolution frequency). Detailed pseudocode.]}

\subsubsection{Experimental Design}\label{sec:experimental_design}

\textit{[To be written. Three experimental settings, each with three trials:]}

\begin{enumerate}
    \item \textit{Three scaffold tiers: Baseline (simple) / Baseline + AutoEvolve / Fully Custom PokeAgent}
    \begin{itemize}
        \item \textit{Test Gemini 3.1-pro-preview}
        \item \textit{Test Gemini 3-flash-preview}
    \end{itemize}
    \item \textit{Sophisticated Supervisor: powerful model for evolution, weaker model for orchestrator}
    \item \textit{Bootstrapped Transfer Learning: equip weaker orchestrator with auto-evolved harness from stronger model + supervisor}
\end{enumerate}

\subsubsection{Metrics}\label{sec:autoevolve_metrics}

\textit{[To be written. Standard speedrunning evaluation criteria; qualitative observations for each setting.]}

\subsection{Results}\label{sec:results}

\subsubsection{Three Settings}\label{sec:three_settings_results}

\textit{[To be written. Speedrunning evaluation metrics and qualitative results for baseline / baseline+AutoEvolve / PokeAgent.]}

\subsubsection{Sophisticated Supervisor}\label{sec:supervisor_results}

\textit{[To be written. Speedrunning evaluation metrics and qualitative results for supervisor configuration.]}

\subsubsection{Bootstrapped Transfer Learning}\label{sec:transfer_results}

\textit{[To be written. Speedrunning evaluation metrics and qualitative results for transfer configuration.]}

\subsection{Discussion}\label{sec:autoevolve_discussion}

\textit{[To be written. How did AutoEvolve compare to simplest baseline and fully custom PokeAgent? Main failure/success cases. Implications for extending to other domains. Literature-supported analysis.]}
